% 원제: 1. 개요
\chapter{개요}

\section{과제 목적과 연구 질문}

800 V급 배터리는 같은 전력을 더 작은 전류로 전달해 급속충전과 배선 손실에 유리하다\cite{b1}. 그러나 차량 안팎의 전압은 하나로 통일되지 않는다. 차량은 다음 기능을 함께 수행해야 한다.

\begin{itemize}
\item 400·500 V급 충전설비와의 호환
\item 400 V급 보조부하(컴프레서, 히터)에 대한 전원 공급
\item 48 V·12 V 전원망의 절연 변환과 고장 시 공급 연속성
\item 여러 부하가 겹칠 때 변환기 정격을 넘지 않도록 하는 전력 관리
\end{itemize}

전전력 절연 변환기는 정밀 제어가 쉽지만 150 kW 전부를 처리하는 자기소자, 반도체와 냉각이 필요하다\cite{b2}. 2:1 고정비 변환기는 효율이 높지만\cite{b3} 출력 전압이 배터리 전압을 따라 움직인다. PRISM-Hub는 대전력 전달과 조정을 분리해 이 상충을 다룬다. 고정비 경로와 소용량 부분전력 조정단의 결합, 직렬 전압보상은 선행연구에 원리가 있다\cite{b4,b5,b6}. 본 과제는 이 원리를 차량의 충전, 400 V 보조버스와 48 V·12 V 전원망에 배치하고, 포트 정격·운전 모드·고장 강등을 하나의 시스템으로 검증하는 데 범위를 둔다. 연구 질문은 다음 세 가지다.

\begin{itemize}
\item Q1(정격 분리). 150 kW급 충전전력을 조정단에서 분리한 상태에서, 소용량 조정단이 보조버스의 차동전력만으로 배터리 전 범위를 조정할 수 있는가.
\item Q2(강건 제어와 부하 관리). 하나의 제어기와 인과적 부하 승인 계층이 부품 공차, 디지털 제어 지연, 부하 중첩과 고장 조건에서 포트 정격과 전압 품질을 함께 지킬 수 있는가.
\item Q3(기준 구조 대비 이점). 이 구조가 조정 없는 고정비 버스나 전전력 변환기를 쓰는 구조보다 손실, 정격 또는 부품 수에서 측정 가능한 이점을 주는가. 준다면 어떤 팩 등급과 부품 등급에서인가.
\end{itemize}

세 질문은 다섯 가설(H1–H5)과 수락 기준으로 판정한다(3.2절). Q3의 답이 사양 팩에서 부정적으로 나와(6장) 구조를 개선했으므로, Q3에는 적용 범위(7장)와 산업에서 쓰는 방식 대비 성립 조건(13장)을 묻는 부분을 더했다. 이번 단계의 범위는 MATLAB/Simulink와 Python을 이용한 시뮬레이션 기반 설계 검증이며, 하드웨어 실증은 후속 단계에서 수행한다.

\section{보고 범위와 기존 보고 대비 변경 사항}

본 보고서는 PRISM-Hub의 설계·검증 결과 전체를 다룬다. 2026년 10월 15일 보고까지의 결과는 5장에 요약하고, 10월 25일 보고의 결과(손실·정격 비교와 R2 정의, 창 조정 모듈, 충전 통합, 보호와 고장 강등)는 6–9장에 싣는다. 이번 보고에서 추가한 검증은 10–13장이고, 향후 계획은 16장에 적는다. 10월 25일 보고 대비 변경 사항은 표 \ref{t:changes}와 같다.

{\footnotesize
\begin{longtable}{L{\dimexpr 0.2069\linewidth-2\tabcolsep\relax}L{\dimexpr 0.2759\linewidth-2\tabcolsep\relax}L{\dimexpr 0.4310\linewidth-2\tabcolsep\relax}L{\dimexpr 0.0862\linewidth-2\tabcolsep\relax}}
\caption{10월 25일 보고 대비 주요 변경 사항}\label{t:changes}\\
\toprule
\tabhead \textbf{항목} & \textbf{10월 25일 보고} & \textbf{본 보고} & \textbf{관련 절} \\
\midrule
\endfirsthead
\toprule
\tabhead \textbf{항목} & \textbf{10월 25일 보고} & \textbf{본 보고} & \textbf{관련 절} \\
\midrule
\endhead
\bottomrule
\endlastfoot
조정형 Stage C & 미설계(고정이득 3 kW 후보만 검증) & 동기 벅 + CLLC 직류변압기 4.8:1, 5 kW, 입력 270–504 V. 평균모델 8/8, 회로 TB14 → TB15 5/5(E22) & 2.6, 10.1–10.2절 \\
R2 통합 회로 & Stage A·Stage C를 이상 전원·정전력 등가로 두고 확인(TB08–TB13) & 실제 Stage A·조정형 Stage C 회로로 기본형 3/3(TB16), 창 조정 모듈 강등 중 48 V 2/2(TB17)(E23) & 10.3절 \\
Stage A 경부하 운전 & 2위상 영교차 추종 & 경부하 위상 셰딩(멈춤 15 A, 복귀 22 A·순시 36 A). 정상 26/26, 과도 9/10, 회로 5/5·2/2(E24, TB18·TB19) & 2.6절, 11장 \\
충전 중 Stage A 정지 & 남은 과제(충전기 전류 차단과의 협조) & 정지 순서 3종, V\textsubscript{A} 클램프 요구(468/463 V, 1 Ω, 7 kJ), 회로 3/4(E26, TB20) & 2.6절, 12장 \\
기준 구조 대비 이점(Q3) & 전용 승압기·전전력 조정기 대비(E15, E18) & 모터·인버터 승압 충전, 800 V급 보조부품 설계 대비(E25). 선행연구 12건과의 비교, 기여와 성립 조건 & 13장 \\
설계 요구 & 게이트 분해능 ≤ 0.625 ns, 충전전류 증가 램프, DAB 정격 1.8배·첨두 상한 등 & + Stage C 데드타임 창, 위상 셰딩 문턱, V\textsubscript{A} 클램프, Stage A 위상당 고정 손실 증가 < 14.8 W(공심 공진 인덕터) & 3.1절 \\
소자 수준 회로 모델 & TB01–TB13 & TB01–TB20(조정형 Stage C, R2 통합, 위상 셰딩, 충전 중 정지) & 14장 \\
\end{longtable}
}
